\documentclass[a4paper,12pt]{article}
\usepackage[utf8]{inputenc}
\usepackage[T1]{fontenc}
\usepackage[ngerman]{babel}
\usepackage{geometry}
\usepackage{hyperref}
\usepackage{booktabs}
\geometry{margin=2.5cm}

\begin{document}

\section*{A Comprehensive and Fair Comparison of Two Neural Operators
(with Practical Extensions) based on FAIR Data}

\noindent\textbf{Autoren:} Lu Lu, Xuhui Meng, Shengze Cai, Zhiping Mao,
Somdatta Goswami, Zhongqiang Zhang, George Em Karniadakis \\
\textbf{Journal:} Computer Methods in Applied Mechanics and Engineering,
393, 114778 \\
\textbf{Jahr:} 2022

\subsection*{Zusammenfassung}

Diese Arbeit pr\"asentiert einen methodisch rigorosen und systematischen
Vergleich der zwei zentralen neuronalen Operatoren \emph{DeepONet} und
\emph{Fourier Neural Operator} (FNO) anhand von 16 sorgf\"altig
ausgew\"ahlten Benchmark-Problemen aus der wissenschaftlichen Simulation.
Beide Architekturen werden unter kontrollierten und vergleichbaren
Bedingungen getestet, wobei die Autoren explizit auf einen fairen Vergleich
achten -- ein Aspekt, der in der fr\"uheren Literatur oft vernachl\"assigt
wurde.

\paragraph{DeepONet} basiert auf dem universellen Approximationssatz f\"ur
Operatoren von Chen \& Chen (1995) und besteht aus zwei Teilnetzen: einem
\emph{Branch-Netz}, das die Eingabefunktion verarbeitet, und einem
\emph{Trunk-Netz}, das die Ausgabekoordinaten als Eingabe erh\"alt. Die
Ausgabe ergibt sich als Skalarprodukt beider Teilnetze. Ein wesentlicher
Vorteil von DeepONet ist seine Flexibilit\"at: Das Branch-Netz kann
beliebige Architekturen verwenden (FNN, CNN, ResNet, GNN), und die
Ausgabe kann an beliebigen Punkten im Raum ausgewertet werden, ohne die
Ausgabefunktion diskretisieren zu m\"ussen.

\paragraph{FNO} parametrisiert den Integraloperator direkt im
Fourier-Raum und verarbeitet die Eingabefunktion durch aufeinanderfolgende
Fourier-Schichten, die spektrale Gewichte mit einem lokalen residualen
Pfad kombinieren. FNO ist auf \"aquidistante Gitter und gleiche
Eingabe-/Ausgabe-Dom\"anen beschr\"ankt, bietet daf\"ur aber bei solchen
einfachen Einstellungen h\"aufig eine sehr effiziente Berechnung.

\paragraph{Neue Erweiterungen} Die Autoren entwickeln mehrere praxisrelevante
Erweiterungen. F\"ur DeepONet werden unter anderem vorberechnete
POD-Moden (Proper Orthogonal Decomposition) als Trunk-Netz eingesetzt
(\emph{POD-DeepONet}), harte Randbedingungen durch Modifikation der
Netzwerkausgabe durchgesetzt sowie eine Reskalierung basierend auf
Varianzanalyse eingef\"uhrt. F\"ur FNO werden Erweiterungen auf
unterschiedliche Dom\"anen (\emph{dFNO+}) und komplexe Geometrien
(\emph{gFNO+}) vorgeschlagen. Zus\"atzlich wird f\"ur beide Architekturen
eine Ausgabenormalisierung als wichtige Ma\ss{}nahme zur
Genauigkeitsverbesserung identifiziert.

\paragraph{Wichtige Ergebnisse} Die 16 Benchmarks umfassen Probleme mit
unterschiedlichen Charakteristika: glatte und diskontinuierliche L\"osungen,
einfache und komplexe Geometrien, station\"are und instation\"are Str\"omungen
sowie verrauschte Daten. Die wesentlichen Erkenntnisse lassen sich wie
folgt zusammenfassen:

\begin{itemize}
  \item F\"ur einfache Einstellungen auf \"aquidistanten Gittern zeigen
    DeepONet und FNO vergleichbare Genauigkeit (z.\,B.\ Burgers-Gleichung:
    FNO $1{,}93\,\%$, POD-DeepONet $1{,}94\,\%$ relativer $L^2$-Fehler).

  \item Bei komplexen Geometrien (Pentagramm, Dreieck mit Kerbe) versagt
    der standard FNO vollst\"andig, w\"ahrend DeepONet und POD-DeepONet
    robuste Ergebnisse liefern (Dreieck: POD-DeepONet $0{,}18\,\%$
    gegenüber dgFNO+ $1{,}00\,\%$).

  \item FNO zeigt systematische Schwierigkeiten bei diskontinuierlichen
    L\"osungen, da die Fourier-Approximation dort ihre Schw\"achen hat
    (Gibbs-Ph\"anomen). DeepONet ist hier deutlich robuster.

  \item POD-DeepONet \"ubertrifft in den meisten Szenarien sowohl
    vanilla DeepONet als auch FNO und erreicht die beste Genauigkeit
    bei geringstem Speicherbedarf.

  \item Feature-Erweiterung -- das Hinzuf\"ugen dom\"anenspezifischen
    Vorwissens als zus\"atzliche Eingaben -- verbessert die Leistung
    beider Architekturen signifikant.

  \item Die Ausgabenormalisierung ist eine einfache, aber wirkungsvolle
    Ma\ss{}nahme, die f\"ur beide Architekturen konsequent empfohlen wird.
\end{itemize}

\paragraph{Theoretischer Vergleich} Die Autoren zeigen formal, dass FNO
in seiner kontinuierlichen Form als Spezialfall von DeepONet aufgefasst
werden kann, bei dem das Branch-Netz eine trigonometrische Basis verwendet.
F\"ur die Burgers-Gleichung wird nachgewiesen, dass beide Architekturen bei
gleicher Netzwerkgr\"o\ss{}e $\mathcal{O}(m^3 \ln m)$ eine Genauigkeit von
$\mathcal{O}(m^{-1})$ erreichen k\"onnen.

\subsection*{Bezug zur Masterarbeit}

Diese Publikation ist aus mehreren Gr\"unden direkt relevant f\"ur die
Masterarbeit \emph{Physics-Based Machine Learning for Predicting Flow
Fields Around Settling Crystals}.

\paragraph{Direkter Architekturvergleich} Die Masterarbeit vergleicht
U-Net und FNO als Surrogatmodelle f\"ur Stokes-Str\"omungen. Lu et al.\ liefern
eine methodische Vorlage f\"ur einen fairen Vergleich neuronaler Operatoren
unter kontrollierten Bedingungen, die direkt \"ubertragbar ist. Insbesondere
die Empfehlung, beide Architekturen mit vergleichbaren Parameteranzahlen
und identischen Trainingsdaten zu evaluieren, sollte in der Masterarbeit
explizit aufgegriffen werden.

\paragraph{Robustheit bei geometrischer Komplexit\"at} Das zentrale Ergebnis,
dass FNO bei komplexen oder unregelm\"a\ss{}igen Geometrien deutlich schlechter
abschneidet als DeepONet, ist f\"ur Kristallkonfigurationen mit mehreren
Partikeln hochrelevant. Mehrere Kristalle erzeugen eine intrinsisch
unregelm\"a\ss{}ige r\"aumliche Struktur, bei der FNO aufgrund seiner
Abh\"angigkeit von \"aquidistanten Gittern systematisch benachteiligt sein
k\"onnte. Die Ergebnisse bei U-Net sollten in diesem Kontext interpretiert
werden, da U-Net lokal im physikalischen Raum operiert und damit strukturell
flexibler ist als FNO.

\paragraph{POD als Vorbild f\"ur Residuall\"arnen} Der Einsatz von POD-Moden
als physikalisch motiviertes Trunk-Netz in POD-DeepONet entspricht
konzeptuell dem Residuallernansatz der Masterarbeit ($v_\mathrm{total} =
v_\mathrm{Stokes,analytisch} + \Delta v_\mathrm{gelernt}$): In beiden
F\"allen werden bekannte Strukturen der L\"osung vorab extrahiert, und das
Netzwerk lernt lediglich die Korrektur bzw.\ die Koeffizienten bez\"uglich
einer gut gew\"ahlten Basis. Die starke Leistung von POD-DeepONet
best\"atigt indirekt, dass solche physikalisch motivierten Induktionsbias-
Strategien generell vorteilhaft sind.

\paragraph{Normalisierung und Skalierung} Die in der Arbeit betonte
Bedeutung der Ausgabenormalisierung deckt sich mit den Herausforderungen
der Masterarbeit bei der Normalisierung \"uber verschiedene
Kristallkonfigurationen hinweg. Lu et al.\ zeigen, dass fehlerhafte oder
fehlende Normalisierung die Leistung beider Architekturen stark
verschlechtern kann -- ein Befund, der die in der Masterarbeit entwickelten
physikalisch motivierten Normalisierungsstrategien (per-sample vs.\
physikalisch skalierte Normalisierung) unterst\"utzt.

\paragraph{Diskontinuit\"aten und Grenzschichteffekte} Die
Beobachtung, dass FNO bei diskontinuierlichen L\"osungen und steilen
Gradienten (z.\,B.\ Grenzschichten im Elektrokonvektionsproblem)
systematisch schw\"acher abschneidet, ist f\"ur Stokes-Str\"omungen um
Kristalle relevant, da die Geschwindigkeitsfelder starke lokale Gradienten
in der N\"ahe der Kristalloberfl\"achen aufweisen. Hier hat U-Net aufgrund
seiner lokal operierenden Faltungsoperationen m\"oglicherweise strukturelle
Vorteile, w\"ahrend FNO globale spektrale Muster bevorzugt.

\paragraph{Generalisierung auf ungesehene Konfigurationen} Das in der
Masterarbeit zentrale Generalisierungsproblem -- vom Training auf
$N$-Kristall-Konfigurationen zur Vorhersage f\"ur $N{+}m$-Kristalle --
wird von Lu et al.\ nicht direkt adressiert, aber deren Ergebnis, dass FNO
eine gr\"o\ss{}ere Generalisierungsl\"ucke zwischen Training- und Testfehler
aufweist als DeepONet, legt nahe, dass FNO bei starker Verteilungs-
verschiebung (z.\,B.\ neue Kristallanzahl) schlechter generalisieren k\"onnte.

\end{document}